\cgeentry{cge_calibrate}{Calibrate the model to the SAM}

\cgesec{Description}

\cmd{cge_calibrate} computes the base value of every variable and every
parameter of the model from the SAM mapped by \manref{cge_map} and the
elasticities given, sets the default closure, and checks that the base solves
every equation of the model (benchmark replication). It stops with an error if
replication fails.

Elasticities have no defaults: they are part of the specification of the model
and must be stated. \manref{cge_elasticities} gives published values, with
their sources, as values for the options of \cmd{cge_calibrate}.

\cgesec{Quick start}

Calibrate with one elasticity of each kind for every sector
\begin{stcode}
cge_calibrate, sigmava(0.8) sigmaq(2) sigmat(2)
\end{stcode}
As above, with one value of \cmd{sigmava()} per activity and of \cmd{sigmaq()}
and \cmd{sigmat()} per commodity, in the order of \manref{cge_map}
\begin{stcode}
cge_calibrate, sigmava(0.5 1 1.5) sigmaq(0.6 3 1.2) sigmat(0.7 4 1)
\end{stcode}
Linear expenditure system demand, with income elasticities by commodity and a
Frisch parameter
\begin{stcode}
cge_calibrate, sigmava(0.8) sigmaq(2) sigmat(2) incelas(0.6 1 1.4) frisch(-2)
\end{stcode}
A labour force for \cmd{lab}, from a base unemployment rate of 15 per cent,
with a wage curve
\begin{stcode}
cge_calibrate, sigmava(0.8) sigmaq(2) sigmat(2) unemployment(lab:0.15) wagecurve(lab:0.1)
\end{stcode}
Armington elasticities from the elasticity library
\begin{stcode}
cge_elasticities sigmaq, country(TGO) values
cge_calibrate, sigmava(0.8) sigmaq(`r(values)') sigmat(2)
\end{stcode}

\cgesec{Syntax}

\begin{stsyntax}
\cmdb{cge_calibrate}\texttt{,} \optarg{}{sigmava}{numlist}
\optarg{}{sigmaq}{numlist} \optarg{}{sigmat}{numlist} [\argu{options}]
\end{stsyntax}

\begin{opttable}
\optgroup{Main}
\optarg{}{sigmava}{numlist} & substitution between factors in value added, one value or one per activity; required\\
\optarg{}{sigmaq}{numlist} & substitution between imports and domestic goods (Armington), one value or one per commodity; required\\
\optarg{}{sigmat}{numlist} & transformation between exports and domestic sales (CET), one value or one per commodity; required\\
\optarg{fq}{uantities}{matname} & base factor quantities, one row per factor and one column per activity\\
\optgroup{Production}
\optarg{}{sigmaa}{numlist} & substitution between value added and intermediate input, one value or one per activity\\
\optarg{nest}{s}{spec} & nests of value added and their members\\
\optarg{sigman}{est}{spec} & substitution within the nests\\
\optarg{}{sigmaac}{spec} & substitution between the outputs of a commodity from several activities\\
\optgroup{Household demand}
\optarg{inc}{elas}{numlist} & income elasticities, one value or one per commodity\\
\optarg{fris}{ch}{numlist} & Frisch parameter, one value or one per household\\
\optgroup{Fiscal instruments}
\optarg{vatc}{redit}{spec} & share of the VAT on its inputs that each activity recovers\\
\optarg{}{debt}{spec} & public debt at home and abroad at the start of the base year\\
\optarg{int}{erest}{spec} & interest rates on the debt at home and abroad\\
\optarg{debth}{olders}{spec} & shares of domestic interest received by households and enterprises\\
\optarg{for}{eignborrowing}{\#} & share of the deficit borrowed abroad in dynamic runs\\
\optgroup{Trade}
\optarg{exporto}{nly}{namelist} & commodities whose domestic output is all exported\\
\optarg{exportd}{emand}{spec} & elasticities of export demand\\
\optgroup{Labour market}
\optarg{unemp}{loyment}{spec} & base unemployment rates by factor\\
\optarg{wagec}{urve}{spec} & wage-curve elasticities by factor\\
\optarg{minw}{age}{spec} & minimum real wages, as shares of the base real wage\\
\optarg{smooth}{ing}{\#} & smoothing parameter of the minimum-wage condition; default \texttt{1e-10}\\
\optgroup{Health}
\optarg{}{healthcom}{namelist} & health commodities\\
\optarg{}{healthfac}{namelist} & factors whose efficiency depends on health services\\
\optarg{}{healthelas}{numlist} & elasticities of factor efficiency with respect to health services\\
\optarg{sickd}{ays}{spec} & base shares of working time lost to illness by factor\\
\optarg{sicke}{las}{spec} & elasticities of the share of time lost\\
\optarg{mort}{ality}{\#} & base crude death rate\\
\optarg{morte}{las}{\#} & elasticity of the death rate\\
\end{opttable}

\cgesec{Options}

\cgesubsec{Main}

\begin{optdesc}
\item[\optarg{}{sigmava}{numlist}] gives the elasticity of substitution between
factors in value added: one value for every activity, or one value per activity
in the order of \cmd{activities()} of \manref{cge_map}. A value of 1 selects
the Cobb--Douglas form. Values must be positive.

\item[\optarg{}{sigmaq}{numlist}] gives the elasticity of substitution between
imports and domestic goods in composite supply (the Armington elasticity), one
value or one per commodity, in the order of \cmd{commodities()}. Values must be
positive, and different from 1 for every commodity that has both imports and
domestic sales: the model has no Cobb--Douglas form of the Armington function.
The values of commodities without imports, or without domestic sales, are not
used.

\item[\optarg{}{sigmat}{numlist}] gives the elasticity of transformation between
exports and domestic sales (CET), one value or one per commodity. Values must
be positive. For a commodity whose exports are at least its domestic output
(re-exports, or exports supplied from inventories), exports are drawn from
composite supply instead, and its \cmd{sigmat()} value is the price elasticity
of those exports in equation \eqref{eq:T9}; calibration lists such commodities. The
values of commodities without exports, and of commodities in
\cmd{exportonly()}, are not used.

\item[\optarg{fq}{uantities}{matname}] supplies base factor quantities, for
example employment, as a Stata matrix with one row per factor and one column
per activity, in the order of \manref{cge_map}. Quantities must be positive
exactly where the SAM has factor payments. The base price of each factor is
then its total payments over its total quantity, and the wage distortion
factors make each payment equal to price times quantity. Without
\cmd{fquantities()}, factor prices are 1 in the base and quantities equal
payments.
\end{optdesc}

\cgesubsec{Production}

\begin{optdesc}
\item[\optarg{}{sigmaa}{numlist}] gives the elasticity of substitution between
value added and aggregate intermediate input, the top nest of production, one
value or one per activity: 0 is the Leontief nest, 1 the Cobb--Douglas form,
and other positive values a CES nest, under which an activity uses relatively
less of whichever input becomes dearer. Without \cmd{sigmaa()}, the top nest is
Leontief in every activity. Activities without intermediate inputs keep the
Leontief nest, and are listed in a note.

\item[\optarg{nest}{s}{spec}] defines nests of value added, separated by
commas, each a name followed by a colon and its members, factors or
commodities, for example \cmd{nests(labour: lab_hi lab_lo, ke: cap com_ene)}.
Each nest is a CES or Cobb--Douglas aggregate of its members, and value added
combines the factors in no nest with the nests, with the elasticity
\cmd{sigmava()}. A commodity in a nest, such as energy, leaves intermediate
input: activities substitute between it and the other members of its nest. A
factor or commodity may be in one nest only, nests do not contain other nests,
and a nest's name may not be that of a factor, a commodity or another nest.
Results list the quantity (\cmd{QNA}) and price (\cmd{PNA}) of each nest in
each activity (\secref{sec:nests}).

\item[\optarg{sigman}{est}{spec}] gives the elasticities of substitution within
the nests, one positive value for every nest or
\argu{nest}\cmd{:}\argu{value} pairs for each. A value of 1 selects the
Cobb--Douglas form. \cmd{nests()} and \cmd{sigmanest()} must be given together.

\item[\optarg{}{sigmaac}{spec}] combines the outputs of a commodity from
several activities by a CES function, instead of treating them as perfect
substitutes, so that each activity is paid its own price for the commodity
(\cmd{PXAC}; \secref{sec:outputces}). Give \argu{commodity}\cmd{:}\argu{value}
pairs, for example \cmd{sigmaac(com_man:2)}, for commodities produced by more
than one activity, or one value for every such commodity. Values must be
positive; a value of 1 selects the Cobb--Douglas form.
\end{optdesc}

\cgesubsec{Household demand}

\begin{optdesc}
\item[\optarg{inc}{elas}{numlist} and \optarg{fris}{ch}{numlist}] select linear
expenditure system (LES) demand instead of Cobb--Douglas demand.
\cmd{incelas()} gives positive income elasticities, one value or one per
commodity; they are scaled so that the marginal budget shares of each
household sum to one, which keeps their ratios. \cmd{frisch()} gives the Frisch
parameter, one value or one per household, below $-1$; supernumerary spending
is then $-1/\textit{frisch}$ of consumption spending. \cmd{frisch()} is
required with \cmd{incelas()}; given alone, it sets every income elasticity to
1. Household consumption must be non-negative. Home consumption of the output
of an activity takes the income elasticity of the activity's paired commodity.
Under LES demand, the equivalent variation reported by \manref{cge_results}
uses the LES expenditure function.
\end{optdesc}

\cgesubsec{Fiscal instruments}

\begin{optdesc}
\item[\optarg{vatc}{redit}{spec}] gives, for a model with a VAT account (see
\cmd{vattax()} and \cmd{vatshare()} of \manref{cge_map}), the share of the VAT
on its inputs that each activity recovers as a credit:
\argu{activity}\cmd{:}\argu{share} pairs or one share for every activity,
between 0 and 1. Activities not listed have share 1, full credit, so that VAT
falls on final demand only; a share of 0 is an exempt or unregistered activity,
which bears VAT on its inputs. The VAT rate of each commodity is calibrated so
that the VAT on its uses, net of credits, equals the VAT that the SAM records
(\secref{sec:vat}).

\item[\optarg{}{debt}{spec} and \optarg{int}{erest}{spec}] give public debt at
the start of the base year and its interest rates, as \cmd{domestic:}\argu{\#}
and \cmd{foreign:}\argu{\#}; a part not given is zero. Domestic debt is in SAM
units, foreign debt in foreign currency (equal to SAM units at the base
exchange rate). Debt must be non-negative and interest rates above $-1$.
Interest is taken to be recorded within the government's transfers to
domestic institutions and abroad; it is separated from them and then paid in
nominal terms on the debt, so that the base is unchanged and interest rates
and debt stocks can be shocked. A transfer smaller than the interest it must
contain is an error. In dynamic runs (\manref{cge_dynamic}) the deficit adds to
the debt (\secref{sec:debt}).

\item[\optarg{debth}{olders}{spec}] gives the shares of domestic interest
received by households and enterprises, as \argu{institution}\cmd{:}\argu{share}
pairs that sum to one. By default the shares are proportional to the
government's transfers to each.

\item[\optarg{for}{eignborrowing}{\#}] is the share of the deficit borrowed
abroad in dynamic runs, between 0 and 1. By default it is the foreign share of
the base debt.
\end{optdesc}

\cgesubsec{Trade}

\begin{optdesc}
\item[\optarg{exporto}{nly}{namelist}] lists commodities whose domestic output
is all exported: their exports, net of export taxes and export margins, must
equal their domestic output. Their output is then exported directly, at the
export price, without domestic sales, and may pay export taxes
(\secref{sec:exportonly}). Without this option such a commodity's exports are
drawn from composite supply; calibration lists, in a note, the commodities whose
exports equal domestic output.

\item[\optarg{exportd}{emand}{spec}] gives elasticities of export demand, as
\argu{commodity}\cmd{:}\argu{elasticity} pairs or one value for every
commodity exported from domestic output. The world export price of these
commodities then falls as more is exported, and a shock to \cmd{pwe} shifts
world demand (\secref{sec:exportdemand}). Without \cmd{exportdemand()}, the
country is small in world markets: it exports any quantity at a fixed world
price.
\end{optdesc}

\cgesubsec{Labour market}

\begin{optdesc}
\item[\optarg{unemp}{loyment}{spec}] gives base unemployment rates, as
\argu{factor}\cmd{:}\argu{rate} pairs with rates at least 0 and below 1, for
example \cmd{unemployment(lab:0.15)}. SAMs record employment, not
unemployment, so the rate comes from a labour force survey. The factor's labour
force is then its base employment divided by one minus the rate; reports show
the labour force and the unemployment rate, and the closures \cmd{wagecurve}
and \cmd{minwage} of \manref{cge_closure} become available. A factor with
unemployment in the base has a minimum wage at its base real wage
(\chapref{chap:labour}).

\item[\optarg{wagec}{urve}{spec}] gives positive wage-curve elasticities, as
\argu{factor}\cmd{:}\argu{elasticity} pairs, for factors with unemployment in
the base: under the closure \cmd{wagecurve} the real wage falls by about
\argu{elasticity} per cent for each per cent rise in the unemployment rate
(Blanchflower and Oswald 1994). A small elasticity approaches a fixed real
wage, a large one full employment.

\item[\optarg{minw}{age}{spec}] gives minimum real wages as shares of the base
real wage, \argu{factor}\cmd{:}\argu{share}. With unemployment in the base the
floor is the base real wage, a share of 1, the only value allowed. With full
employment in the base (no rate, or a rate of 0, in \cmd{unemployment()}), give
a share below 1, for example \cmd{minwage(lab:0.95)}: the floor then binds only
after the real wage falls by 5 per cent.

\item[\optarg{smooth}{ing}{\#}] sets the smoothing parameter $\mu$ of the
minimum-wage condition; the default is \texttt{1e-10}, and it must be positive.
Under the closure \cmd{minwage}, the condition that either unemployment is zero
or the real wage is at its floor is solved, with the rest of the model, in the
smoothed Fischer--Burmeister form (Fischer 1992; Kanzow 1996), in which the
product of the unemployment rate and the proportional excess of the real wage
over the floor is $\mu$.
\end{optdesc}

\cgesubsec{Health}

\begin{optdesc}
\item[\optarg{}{healthcom}{namelist}, \optarg{}{healthfac}{namelist} and
\optarg{}{healthelas}{numlist}] link factor efficiency to health services.
\cmd{healthcom()} lists the health commodities; \cmd{healthfac()} lists the
factors whose efficiency depends on them, all factors by default; and
\cmd{healthelas()} gives the elasticity of each factor's efficiency with
respect to the volume of health services consumed by households, the
government and the health insurance fund, one value or one per factor of
\cmd{healthfac()}, zero or positive. With elasticity $e$, the factor's
efficiency is $(\V{HS}/\V{HS}^0)^e$, where $\V{HS}$ is that volume and
$\V{HS}^0$ its base value; the factor then enters value added in efficiency
units, while its quantity, supply and wage stay in physical units. Without
these options, or with elasticity 0, efficiency changes only through the
parameter \cmd{af} of \manref{cge_shock}. \cmd{healthelas()} needs
\cmd{healthcom()}, and \cmd{healthfac()} needs \cmd{healthelas()}. The health
commodities must have some consumption in the base. These options must be
typed in full.

\item[\optarg{sickd}{ays}{spec}] gives the base share of working time lost to
illness, as \argu{factor}\cmd{:}\argu{share} pairs with shares at least 0 and
below 1, for example \cmd{sickdays(lab:0.05)} from a labour force or household
survey. Days worked per worker then scale the factor's effective input in
production; employment, wages and unemployment stay per worker
(\secref{sec:sickdays}).

\item[\optarg{sicke}{las}{spec}] gives the elasticity of the share of time lost
with respect to health services per head, one value for every factor in
\cmd{sickdays()} or \argu{factor}\cmd{:}\argu{elasticity} pairs for factors in
\cmd{sickdays()}. With elasticity $\epsilon$, the share lost is $s\,R^\epsilon$,
where $s$ is the share given in \cmd{sickdays()} and $R$ is health services per
head relative to the base.

\item[\optarg{mort}{ality}{\#}] gives the base crude death rate, deaths per head
per period, at least 0 and below 1 (for example 0.008 for 8 per 1,000 a year),
and \optarg{morte}{las}{\#} its elasticity with respect to health services per
head; \cmd{mortelas()} needs \cmd{mortality()}. The death rate is reported; it
does not change the solution of a period, but in dynamic runs it changes the
growth of the population and, with \cmd{deathfactors()} of \manref{cge_dynamic},
of labour from the next period (\secref{sec:mortality}).
\end{optdesc}

Health outcomes with a non-zero elasticity need \cmd{healthcom()}. Their
elasticities are hypotheses about how health services become health and
should come from studies of the country or region; the package supplies none.

\cgesec{Remarks and examples}

Calibration chooses the parameters of the model so that the SAM is a solution
of it, the base. \cmd{cge_calibrate}
\begin{enumerate}
\item sets the base prices of activities, commodities and value added, the
exchange rate and world import prices to one, so that base quantities are
values of the SAM, and sets factor prices to one unless \cmd{fquantities()}
gives factor quantities;
\item computes tax rates, shares and the other parameters from the SAM and the
elasticities;
\item checks the identities that a balanced SAM implies, for example that
composite supply equals the demand for each commodity, and stops with error
459 if one fails;
\item sets the default closure: a flexible exchange rate with fixed foreign
savings, flexible government savings, savings-driven investment, the consumer
price index as numeraire, and mobile, fully employed factors
(\manref{cge_closure});
\item checks benchmark replication: every equation of the model must hold at
the base, with a largest scaled residual of at most $10^{-9}$, or the command
stops with error 9. User equations need not hold at the base, where they act
as shocks; from version 0.7.1 replication leaves them out, and a note says
when they do not hold.
\end{enumerate}
Calibrating again replaces the calibration: stored simulations, shocks,
constraints and closure swaps are discarded and the default closure is set.
User variables and equations (\manref{cge_variable}, \manref{cge_equation}) are
kept and compiled again from the new base, and user sets and parameters
(\manref{cge_set}, \manref{cge_parameter}) are kept, the parameters at their
base values.

Each elasticity option takes one value for every activity or commodity, or one
value per activity (\cmd{sigmava()}, \cmd{sigmaa()}) or commodity
(\cmd{sigmaq()}, \cmd{sigmat()}, \cmd{incelas()}) in the order of the lists of
\manref{cge_map}. A failure of benchmark replication points to a problem in the
SAM or in the mapping, not in the solver.

\cgesubsec{Example 1: one elasticity of each kind}

The synthetic SAM is mapped as in Example~1 of \manref{cge_map} and calibrated:
\stlog{cal_uniform}
The model has 93 equations and, under the default closure, as many free
variables. The largest scaled residual at the base is of the order of the
machine precision. The results are stored in \cmd{r()}:
\stlog{cal_return}

\cgesubsec{Example 2: elasticities by sector}

The elasticities may differ by activity and by commodity; a value of 1 for
\cmd{sigmava()} gives Cobb--Douglas value added in manufacturing:
\stlog{cal_sector}
A value of 1 for \cmd{sigmaq()} is not allowed for a commodity with imports
and domestic sales:
\stlog{cal_sigmaq1}

\cgesubsec{Example 3: linear expenditure system demand}

With \cmd{incelas()} and \cmd{frisch()}, households first buy subsistence
quantities and spend the rest of their budget in fixed marginal shares:
\stlog{cal_les}
The model has the same number of equations; only the parameters of household
demand differ.

\cgesubsec{Example 4: exports at least as large as domestic output}

In the supply-use layout of the synthetic SAM, \cmd{com_agr} exports 105
against domestic output of 100, with a fall in its stocks, and \cmd{com_fuel},
which is only imported, re-exports part of its imports:
\stlog{cal_sut}
The exports of these two commodities are drawn from composite supply, by
equation \eqref{eq:T9}, and the \cmd{sigmat()} value is the price elasticity of those
exports.

\cgesubsec{Example 5: an extension}

The options of the extensions add equations and parameters to the standard
model. A base unemployment rate gives labour a labour force, and a wage-curve
elasticity gives the closure \cmd{wagecurve} its slope:
\stlog{cal_labour}
The labour force is base employment, 135, divided by $1 - 0.15$. The model has
two more equations, the wage curve and the minimum-wage condition of
\cmd{lab}; under the default closure, which keeps labour fully employed, both
are inert, and \cmd{cge_closure, factors(lab:wagecurve)} or
\cmd{factors(lab:minwage)} brings one of them into use. The real minimum wage is the
base real wage divided by $1 + \mu/u^0$, where $u^0$ is the base unemployment
rate, so that the smoothed condition holds at the base.

The options of the extensions, and the chapters of Part~II that describe them:

\begin{cgetable}{>{\raggedright\arraybackslash}p{0.30\linewidth}>{\raggedright\arraybackslash}p{0.43\linewidth}>{\raggedright\arraybackslash}p{0.21\linewidth}}
\toprule
Extension & Options & Part II\\
\midrule
top nest of production & \cmd{sigmaa()} & \secref{sec:production}\\
nests of value added & \cmd{nests()}, \cmd{sigmanest()} & \secref{sec:nests}\\
output from several activities & \cmd{sigmaac()} & \secref{sec:outputces}\\
VAT and public debt & \cmd{vatcredit()}, \cmd{debt()}, \cmd{interest()}, \cmd{debtholders()}, \cmd{foreignborrowing()} & \chapref{chap:fiscal}\\
trade & \cmd{exportonly()}, \cmd{exportdemand()} & \chapref{chap:trade}\\
labour market & \cmd{unemployment()}, \cmd{wagecurve()}, \cmd{minwage()}, \cmd{smoothing()} & \chapref{chap:labour}\\
health & \cmd{healthcom()}, \cmd{healthfac()}, \cmd{healthelas()}, \cmd{sickdays()}, \cmd{sickelas()}, \cmd{mortality()}, \cmd{mortelas()} & \chapref{chap:health}\\
\bottomrule
\end{cgetable}

\cgesec{Stored results}

\cmd{cge_calibrate} stores the following in \cmd{r()}:
\begin{storedres}
\resgroup{Scalars}
\res{r(N_equations)} & number of equations\\
\res{r(max_residual)} & largest scaled residual at the base\\
\end{storedres}

\cgesec{Methods and formulas}

Let $\SAM(r, c)$ be the payment from column account $c$ to row account $r$ of
the mapped SAM, after the rewrites of \manref{cge_map}. Base values carry a
superscript 0. The base prices
$\V{PA}^0_a = \V{PX}^0_c = \V{PD}^0_c = \V{PE}^0_c = \V{PQ}^0_c = \V{PVA}^0_a = 1$,
$\V{EXR}^0 = 1$ and $\V{pwm}_c = 1$ (\secref{sec:normalisation}); without
\cmd{fquantities()}, $\V{W}^0_f = \V{WFDIST}^0_{fa} = 1$.
Given base factor quantities $\V{QF}^0_{fa}$,
\begin{equation*}
\V{W}^0_f = \frac{\sum_a \SAM(f, a)}{\sum_a \V{QF}^0_{fa}}, \qquad
\V{WFDIST}^0_{fa} = \frac{\SAM(f, a)}{\V{W}^0_f\,\V{QF}^0_{fa}} .
\end{equation*}

The tax rates are taken from the tax accounts:
\begin{equation*}
tm_c = \frac{\SAM(\textit{imp}, c)}{\SAM(\textit{row}, c)}, \qquad
te_c = \frac{\SAM(\textit{exp}, c)}{\SAM(c, \textit{row})}, \qquad
ta_a = \frac{\SAM(\textit{act}, a)}{\V{QA}^0_a}, \qquad
tq_c = \frac{\SAM(\textit{sal}, c)}{\V{QQ}^0_c},
\end{equation*}
where $\V{QA}^0_a$ is the column total of activity $a$ and, without margins by
market and VAT,
$\V{QQ}^0_c = \V{QD}^0_c + (1 + tm_c)\,\V{QM}^0_c + \SAM(\textit{sal}, c) +
\sum_{t} \SAM(t, c)$, the last sum over the margin accounts.

The elasticities enter the functional forms through the exponents
\begin{equation*}
\rho^{va}_a = \frac{1}{\sigma^{va}_a} - 1, \qquad
\rho^{q}_c = \frac{1}{\sigma^{q}_c} - 1, \qquad
\rho^{t}_c = \frac{1}{\sigma^{t}_c} + 1,
\end{equation*}
and the share parameters follow from the first-order conditions at the base.
For CES value added, with $w_{fa} = \V{W}^0_f\,\V{WFDIST}^0_{fa}$,
\begin{equation*}
\delta^{va}_{fa} = \frac{w_{fa}\,(\V{QF}^0_{fa})^{1+\rho^{va}_a}}
  {\sum_{f'} w_{f'a}\,(\V{QF}^0_{f'a})^{1+\rho^{va}_a}}, \qquad
ad_a = \V{QVA}^0_a \Big/ \Bigl[\sum_f \delta^{va}_{fa}\,(\V{QF}^0_{fa})^{-\rho^{va}_a}\Bigr]^{-1/\rho^{va}_a},
\end{equation*}
and for Cobb--Douglas value added $\delta^{va}_{fa} = \SAM(f, a)/\V{QVA}^0_a$.
The Armington and CET shares are
\begin{equation*}
\delta^{q}_c = \frac{r^q_c}{1 + r^q_c}, \quad
r^q_c = \frac{\V{PM}^0_c}{\V{PD}^0_c}\Bigl(\frac{\V{QM}^0_c}{\V{QD}^0_c}\Bigr)^{1+\rho^{q}_c};
\qquad
\delta^{t}_c = \frac{1}{1 + r^t_c}, \quad
r^t_c = \frac{\V{PD}^0_c}{\V{PE}^0_c}\Bigl(\frac{\V{QE}^0_c}{\V{QD}^0_c}\Bigr)^{\rho^{t}_c - 1},
\end{equation*}
with the complements $1 - \delta$ computed from the same ratios, not by
subtraction, so that no precision is lost when a share is close to one.

Under LES demand, with budget shares $w_{ch} = \SAM(c, h)/\V{EH}^0_h$, income
elasticities $\varepsilon_c$ and Frisch parameters $\phi_h < -1$,
\begin{equation*}
\beta_{ch} = \frac{\varepsilon_c\,w_{ch}}{\sum_{c'} \varepsilon_{c'}\,w_{c'h}}, \qquad
\gamma_{ch} = \V{QH}^0_{ch} + \frac{\beta_{ch}\,\V{EH}^0_h}{\V{PH}^0_c\,\phi_h},
\end{equation*}
so that supernumerary spending is $-\V{EH}^0_h/\phi_h$; under Cobb--Douglas
demand $\beta_{ch} = w_{ch}$ and $\gamma_{ch} = 0$.

Benchmark replication evaluates every equation $i$ of the model at the base,
as $L_i(x^0) = R_i(x^0)$, and requires
\begin{equation*}
\max_i \frac{|L_i(x^0) - R_i(x^0)|}{e_i} \le 10^{-9}, \qquad
e_i = \max\bigl(|L_i(x^0)|, |R_i(x^0)|\bigr),
\end{equation*}
with $e_i = 1$ where that maximum is below $10^{-9}$. The same scales $e_i$ are
used by \manref{cge_solve}.

\Chapref{chap:calibration} gives every calibration formula, including those of
the top nest, margins, VAT, institutions and the health insurance fund, and the
chapters on the extensions give theirs.

\cgesec{References}

\begin{list}{}{\setlength{\leftmargin}{2em}\setlength{\itemindent}{-2em}%
  \setlength{\itemsep}{2pt}\setlength{\parsep}{0pt}}
\item Blanchflower, D. G., and A. J. Oswald. 1994. \textit{The Wage Curve}.
Cambridge, MA: MIT Press.
\item Fischer, A. 1992. A special Newton-type optimization method.
\textit{Optimization} 24: 269--284.
\item Kanzow, C. 1996. Some noninterior continuation methods for linear
complementarity problems. \textit{SIAM Journal on Matrix Analysis and
Applications} 17: 851--868.
\end{list}

\cgesec{Also see}

\begin{alsosee}
\seeentry{cge_map}{Assign SAM accounts to model roles}
\seeentry{cge_elasticities}{Published elasticities, with their sources}
\seeentry{cge_check}{Run model validation checks}
\seeentry{cge_parameters}{The calibrated parameters as a dataset}
\seeentry{cge_closure}{Choose the model closure}
\seeentry{cge_solve}{Solve the model}
\seechapter{chap:calibration}{Calibration}
\end{alsosee}
